\
\documentclass[11pt]{article}
\usepackage[utf8]{inputenc}
\usepackage[T1]{fontenc}
\usepackage[margin=1in]{geometry}
\usepackage{hyperref}
\usepackage{enumitem}
\title{Event-Driven Robot Coding}
\author{Piter Garcia}
\date{March 20, 2026}
\begin{document}
\maketitle
\section*{Quick Scan}
\begin{itemize}[leftmargin=*]
\item Grade: Grade 4
\item Packet focus: robot sensors, event-response logic, and debugging input-output behaviors
\item Class blocks: Day 1 1:15-2:10 Baris, Day 2 1:15-2:10 Fowler, Day 3 1:15-2:10 Schrank, Day 4 1:15-2:10 Autore
\item Reference file: event-driven-robot-coding.bib
\end{itemize}
\section*{School Planning Goals and Inclusion}
\begin{itemize}[leftmargin=*]
\item What is expected: make the CS learning target, proof, and revision path visible.
\item What we achieved: This packet is grounded in local evidence notes and cleaned transcripts from real teaching sessions.
\item What needs improvement: clearer proof, sharper assessment, and fewer hidden assumptions in directions.
\item How to improve: keep one artifact, one observation, and one revision note after reteaching.
\item Inclusion focus: visual modeling, chunked directions, role clarity, and flexible pacing supports.
\end{itemize}
\section*{Official References}
\begin{itemize}[leftmargin=*]
\item \url{https://csteachers.org/k12standards/interactive/}
\item \url{https://dmmedia.sphero.com/email-marketing/Sphero-Edu/SpheroEdu-k12-teacher-resource-guide-v1_updated050818.pdf}
\end{itemize}
\end{document}
